\documentclass[11pt,a4paper]{article}

\usepackage[margin=2.2cm]{geometry}
\usepackage[T1]{fontenc}
\usepackage[utf8]{inputenc}
\usepackage{lmodern}
\usepackage{microtype}
\usepackage{parskip}
\usepackage{booktabs}
\usepackage{tabularx}
\usepackage{array}
\usepackage{enumitem}
\usepackage{graphicx}
\usepackage{hyperref}
\usepackage{xcolor}

\hypersetup{hidelinks}
\setlist[itemize]{leftmargin=*, itemsep=2pt, topsep=2pt}
\setlength{\parskip}{0.5em}
\setlength{\parindent}{0pt}
\graphicspath{{figures/}}

\newcolumntype{Y}{>{\raggedright\arraybackslash}X}

\title{\textbf{SHA-256 Hardware Acceleration in Croc}\\
\large Current Implementation Status Report}
\author{Radu-Constantin Cretu \and Alexander Huwiler}
\date{June 2026}

\begin{document}
\maketitle

\begin{abstract}
This report documents the current implementation status of our Croc-based
system-on-chip extension that adds a dedicated SHA-256 accelerator in the
\texttt{user\_domain}. The project extends the baseline Croc design with new RTL,
software support, directed functional tests, broad end-to-end verification, and
frontend implementation checks. The accelerator exposes a memory-mapped
programming model at \texttt{0x2000\_1000}, supports both single-block and
multi-block chained hashing, and is complemented by the required user ROM at
\texttt{0x2000\_0000} for automatic name readout. At the current project stage,
all SHA-specific tests pass on a fresh Verilator rebuild, and ETH
reference-flow spot checks confirm correct end-to-end hashing and user-ROM
behavior in the intended environment. A fresh Yosys synthesis run on the ETH
reference flow completes successfully, and the OpenROAD floorplan stage now
reaches a valid checkpoint after aligning the top-level pad cells with the
reference I/O library. Placement, routing, final timing closure, and
manufacturing signoff remain the next major tasks.
\end{abstract}

\section{Introduction}
The goal of this project is to add a meaningful hardware improvement to the
reference Croc SoC used throughout the VLSI 2 exercises. We chose SHA-256
acceleration because it is a realistic workload for hardware/software
co-design: the algorithm has a well-defined standard behavior, it is easy to
verify against software reference implementations, and it contains a
computationally intensive compression step that maps naturally to a custom
datapath.

Our implementation does not alter the overall Croc platform structure. Instead,
it adds a new user-domain peripheral that the CPU can access through the
existing OBI-based memory map. This keeps the integration effort aligned with
the course goals: the project is not only about implementing a standalone RTL
block, but about integrating, testing, and eventually implementing a complete
SoC extension under realistic frontend and backend constraints.

\section{System-Level Integration}
The accelerator is integrated into the Croc \texttt{user\_domain}. The user
address map now contains three subordinate targets:
\begin{itemize}
  \item an error subordinate for unmapped addresses,
  \item a user ROM at \texttt{0x2000\_0000},
  \item the SHA-256 accelerator at \texttt{0x2000\_1000}.
\end{itemize}

The integration is implemented in \texttt{rtl/user\_pkg.sv} and
\texttt{rtl/user\_domain.sv}. A dedicated OBI demultiplexer routes requests to
either the ROM or the accelerator. The SHA block raises an interrupt through
\texttt{interrupts\_o[0]}, which appears as \texttt{IRQ\_SHA256\_ACCEL = 20} on
the software side.

The user ROM was added to satisfy the project requirement that the fabricated
chip must be able to output the students' names from the start of the user
domain. The ROM returns the zero-terminated string
\texttt{"Radu-Constantin Cretu, Alexander Huwiler"} in little-endian 32-bit
words, which allows the software infrastructure to read and print it as a C
string.

\begin{figure}[t]
\centering
\includegraphics[width=\textwidth]{Figure1.pdf}
\caption{System-level integration of the user ROM and the SHA-256 accelerator
inside Croc's \texttt{user\_domain}. Requests are routed through the local OBI
infrastructure, while the accelerator exposes a dedicated completion interrupt
path back to the processor-visible interrupt fabric.}
\label{fig:user_domain_integration}
\end{figure}

\section{SHA-256 Accelerator Architecture}
\subsection{Programming Model}
The software-visible programming model is defined in
\texttt{rtl/sha256/sha256\_accel\_pkg.sv},
\texttt{sw/lib/inc/sha256\_accel.h}, and
\texttt{sw/lib/src/sha256\_accel.c}. The register map is intentionally compact:

\begin{itemize}
  \item \texttt{CTRL} at offset \texttt{0x00}
  \item \texttt{STATUS} at offset \texttt{0x04}
  \item \texttt{INFO} at offset \texttt{0x08}
  \item \texttt{BLOCK[0..15]} at offsets \texttt{0x40} to \texttt{0x7C}
  \item \texttt{DIGEST[0..7]} at offsets \texttt{0x80} to \texttt{0x9C}
\end{itemize}

The control register implements three bits:
\begin{itemize}
  \item \texttt{START}: launch a new compression,
  \item \texttt{IRQ\_EN}: enable completion interrupts,
  \item \texttt{CHAIN}: continue hashing from the previously computed internal
  state.
\end{itemize}

The status register reports:
\begin{itemize}
  \item \texttt{READY},
  \item \texttt{BUSY},
  \item \texttt{DONE},
  \item \texttt{DIGEST\_VALID},
  \item \texttt{IRQ\_PENDING},
  \item \texttt{START\_ERR}.
\end{itemize}

This interface allows software to write one fully padded 512-bit block, start a
fresh hash or a chained continuation, wait for completion by polling or by
interrupt, and then read back the resulting 256-bit digest.

\subsection{Datapath and Control Strategy}
The accelerator uses a compact iterative architecture. It is not unrolled and
not pipelined. Each \texttt{START} command performs one SHA-256 compression over
64 rounds, with one round executed per cycle. This keeps the design small and
appropriate for integration into the fixed project area.

The core hardware state consists of:
\begin{itemize}
  \item 16 block input registers for the current 512-bit message block,
  \item 8 digest registers for the externally visible 256-bit result,
  \item 8 chaining state registers that hold the current hash state across
  chained blocks,
  \item 8 working registers \texttt{a} through \texttt{h},
  \item a 6-bit round counter,
  \item a 16-word rolling message-schedule buffer.
\end{itemize}

When software issues a fresh \texttt{START}, the internal state is initialized
to the standard SHA-256 IV. When software issues \texttt{START} together with
\texttt{CHAIN}, the core resumes from the previously stored chaining state.
After round 63, the final working values are added to the current hash state,
the digest registers are updated, and the completion flags are asserted.

Padding is currently performed in software. This means the hardware is
responsible for the compression function and state chaining, while the software
runtime is responsible for turning arbitrary-length messages into one or more
properly padded 512-bit blocks.

\begin{figure}[t]
\centering
\includegraphics[width=\textwidth]{Figure2.pdf}
\caption{Internal block diagram of the SHA-256 accelerator. The design uses a
compact iterative architecture with one round executed per cycle, externally
visible block and digest registers, internal working registers
\texttt{a}--\texttt{h}, and a chaining-state path for multi-block hashing.}
\label{fig:sha_arch}
\end{figure}

\section{Software Support}
To make the accelerator usable from bare-metal software, we added a small
driver library and several dedicated tests in the reference-flow software tree.
The driver provides helpers to:
\begin{itemize}
  \item read and write control and status registers,
  \item write complete message blocks,
  \item start fresh or chained operations,
  \item wait for completion,
  \item read the final digest,
  \item hash one or several blocks through a higher-level API.
\end{itemize}

One driver detail that required special care is the \texttt{START\_ERR} status
handling. A start command issued while the accelerator is already busy raises
\texttt{START\_ERR}, but this should not cause the software wait loop to return
prematurely while the current valid operation is still running. The final driver
implementation therefore waits until either a true completion is observed or a
standalone start error is reported while the accelerator is no longer busy.

\section{Verification Strategy}
\subsection{Directed Functional Tests}
The SHA-specific verification suite is implemented in
\texttt{sw/test/}. The most important directed tests are summarized in
Table~\ref{tab:tests}.

\begin{table}[t]
\centering
\small
\begin{tabularx}{\textwidth}{@{}lYY@{}}
\toprule
Test & Purpose & Current result \\
\midrule
\texttt{test\_sha256\_accel} & Basic one-block hash using the standard padded
\texttt{"abc"} test case & Pass \\
\texttt{test\_sha256\_accel\_api} & Exercise the software driver API for both
single-block and multi-block use & Pass \\
\texttt{test\_sha256\_accel\_multiblock} & Verify explicit chained hashing over
multiple message blocks & Pass \\
\texttt{test\_sha256\_accel\_irq} & Verify interrupt enable, pending, and
completion behavior & Pass \\
\texttt{test\_sha256\_accel\_mmio} & Verify block-register writes, byte enables,
status clearing, and control-register semantics & Pass \\
\texttt{test\_sha256\_accel\_restart} & Verify restart behavior, invalid
restarts while busy, and reset to the SHA-256 IV & Pass \\
\texttt{test\_sha256\_accel\_vectors} & Compare against several known vectors,
including boundary-length cases around SHA padding transitions & Pass \\
\texttt{test\_user\_rom} & Verify correct readout of the required student-name
string from the user ROM & Pass \\
\bottomrule
\end{tabularx}
\caption{Directed verification tests for the new functionality.}
\label{tab:tests}
\end{table}

\subsection{Broad End-to-End Coverage}
In addition to directed tests, we added a broad end-to-end regression
(\texttt{test\_sha256\_accel\_end\_to\_end}) that compares the hardware result
against an in-test software SHA-256 reference implementation. This test builds
properly padded messages in software and then validates the accelerator over a
large set of automatically generated cases.

The current sweep covers:
\begin{itemize}
  \item all message lengths from 0 to 160 bytes for pattern 0,
  \item all message lengths from 0 to 160 bytes for pattern 1,
  \item twelve additional lengths up to 255 bytes for both patterns 0 and 1,
  \item 33 boundary-focused cases for pattern 2, including 55, 56, 57, 63, 64,
  65, 119, 120, 121, 127, 128, 129, 191, 192, 193, and 255 bytes.
\end{itemize}

This corresponds to 379 distinct end-to-end message cases. On the rebuilt local
Verilator setup, the test passes functionally but requires a large timeout
budget because it performs a broad sweep and recomputes a software reference
digest for every case. We therefore reran this test explicitly with an
extended timeout budget and confirmed a clean successful completion on the
current branch.

\subsection{Baseline Regression Status}
We also reran the standard non-SHA regression tests on a freshly rebuilt
Verilator model. The following baseline tests pass on the current branch:
\texttt{helloworld}, \texttt{print\_config}, \texttt{test\_clint},
\texttt{test\_compute}, \texttt{test\_gpio}, \texttt{test\_interrupts},
\texttt{test\_obi\_timer}, \texttt{test\_soc\_ctrl}, and \texttt{test\_uart}.

To reduce the risk of a local-tool-only success, we also reran key tests on the
ETH reference-flow environment. In that setup, \texttt{helloworld},
\texttt{test\_sha256\_accel\_end\_to\_end}, and \texttt{test\_user\_rom}
complete successfully on a fresh Verilator build. These runs are especially
important because they validate the supported SHA configuration in the same
environment that is used for backend implementation.

As a final local sanity pass after rebuilding the software binaries, we reran a
focused short regression consisting of \texttt{helloworld},
\texttt{test\_sha256\_accel}, \texttt{test\_sha256\_accel\_api},
\texttt{test\_sha256\_accel\_irq}, \texttt{test\_sha256\_accel\_mmio},
\texttt{test\_sha256\_accel\_multiblock},
\texttt{test\_sha256\_accel\_restart},
\texttt{test\_sha256\_accel\_vectors}, \texttt{test\_user\_rom},
\texttt{print\_config}, \texttt{test\_soc\_ctrl}, and \texttt{test\_uart}. All
of these tests completed successfully.

\section{Frontend and Synthesis Checks}
On the exact reference-flow \texttt{port-sha} branch, we ran a fresh Yosys
synthesis in the intended ETH 2026 tool environment. The design parses
successfully through \texttt{yosys-slang}, elaboration and process conversion
complete, and the mapped netlist is generated without a fatal synthesis error.
This gives us a valid ETH-side synthesis checkpoint on the correct project tree
before detailed backend optimization begins. The authoritative area and timing
numbers will be taken from the final backend flow once placement, routing, and
signoff checks have stabilized.

We also rebuilt the Verilator model from scratch in the same environment. This
was important because an older stale build directory had been generated against
a different local Verilator installation path. After removing the stale build
artifacts, the simulator rebuilt successfully and all SHA-relevant tests were
rerun on a fresh binary.

\section{Performance Evaluation}
To evaluate performance more defensibly than with a single short-message
example, we benchmarked both the hardware path and a software-only SHA-256
reference over a set of message lengths that exercise important padding and
multi-block boundaries. In this comparison, both paths include software-side
message padding and block preparation; the hardware path then offloads only the
compression and chaining steps to the accelerator. The resulting per-hash cycle
counts are summarized in Table~\ref{tab:perf}. Each reported value is derived
from a 32-iteration simulation run on the rebuilt Verilator model.

To keep this comparison fair, we reran both the hardware benchmark
(\texttt{test\_sha256\_accel\_bench}) and the software-only benchmark
(\texttt{test\_sha256\_sw\_bench}) on the rebuilt simulator with sufficient
timeout budget for the longer software run.

\begin{table}[t]
\centering
\small
\begin{tabularx}{\textwidth}{@{}rrrrr@{}}
\toprule
Length [B] & Blocks & HW cyc/hash & SW cyc/hash & Speedup \\
\midrule
3   & 1 & 1\,379 & 7\,070  & 5.13x \\
55  & 1 & 2\,731 & 8\,422  & 3.08x \\
56  & 2 & 3\,049 & 13\,967 & 4.58x \\
64  & 2 & 3\,257 & 14\,175 & 4.35x \\
119 & 2 & 4\,687 & 15\,605 & 3.33x \\
120 & 3 & 5\,005 & 21\,150 & 4.23x \\
128 & 3 & 5\,213 & 21\,358 & 4.10x \\
255 & 5 & 9\,099 & 35\,698 & 3.92x \\
\bottomrule
\end{tabularx}
\caption{Current benchmark results measured in simulation cycles for selected
message lengths. Both hardware and software paths include software-side padding
and block preparation.}
\label{tab:perf}
\end{table}

These results show a consistent accelerator speedup across all tested message
lengths, with the strongest gain on very short inputs and a still clear
advantage on longer multi-block messages. The transition from 55 to 56 bytes is
particularly useful because it crosses the SHA-256 padding boundary from one
block to two blocks; the expected cost increase is visible in both the hardware
and software measurements. The additional 119-byte and 120-byte points show
the same effect at the next boundary, where the padded message grows from two
blocks to three. Across the currently measured cases, the observed speedup
ranges from 3.08x to 5.13x. Overall, the current implementation provides a real
system-level improvement rather than only a functional hardware addition.

\section{Backend Status and Remaining Work}
At the current project stage, the design has already passed ETH reference-flow
synthesis and the initial OpenROAD floorplan stage. During backend bring-up, we
identified an I/O-library naming mismatch at the top level; after aligning the
RTL pad instances with the \texttt{sg13cmos5l\_*} reference-flow library names,
the floorplan stage now completes and produces a valid backend checkpoint. The
remaining implementation work is:
\begin{itemize}
  \item OpenROAD floorplanning, placement, clock-tree synthesis, routing, and
  finishing,
  \item static timing analysis and collection of final area/timing results,
  \item KLayout GDS generation and final DRC/LVS closure,
  \item post-synthesis and, if feasible, post-layout sanity simulation,
  \item final report consolidation around backend measurements and
  manufacturability results.
\end{itemize}

The authoritative final numbers should be taken from the ETH reference-flow
environment with the correct technology installation, since that is the setup
used for the course backend flow and final submission artifacts.

\section{Conclusion}
So far, the project has progressed well beyond the idea stage. We have:
\begin{itemize}
  \item integrated a SHA-256 accelerator into the Croc user domain,
  \item added the required user ROM for name readout,
  \item implemented the software driver and bare-metal test support,
  \item verified the design through directed tests and a broad end-to-end
  regression,
  \item measured a clear hardware speedup over the software baseline,
  \item established a fresh ETH-side synthesis checkpoint on the reference
  flow,
  \item brought the design through the initial OpenROAD floorplan stage on the
  reference backend setup.
\end{itemize}

The main remaining work is now backend implementation and final physical
verification. From a frontend perspective, the SHA-specific functionality is in
a defensible state for the backend phase to proceed.

\end{document}
